\section{Coefficient-relative resolution}
\label{sec:setup}

All indexing categories in the main construction are finite ordinary posets,
regarded as categories. Their nerves are also used as indexing
infinity-categories when the coefficients are spaces. Write $\B J=|\N J|$
for the classifying space of an ordinary category $J$, with
$\B\varnothing=\varnothing$.

\subsection{Conventions and Kan extensions}

For $k\geq0$ put $\Vk=\Spc_{\leq k}$, and put
$\mathcal V_\infty=\Spc$. These categories admit the limits and colimits
used below. Limits of $k$-types agree with their limits in spaces; colimits
are obtained by taking the colimit in spaces and then truncating. When
groupoids represent 1-types, their limits here are homotopy-coherent
limits, not strict limits in the ordinary category of groupoids. We use
\emph{covariant} diagrams $\Fun(C,\Vk)$. This is the presheaf convention
on $C^{\mathrm{op}}$, not on $C$.

For a full inclusion $u:R\hookrightarrow C$, let $u^*$ be restriction. Its
adjoints are
\begin{equation}\label{eq:triple}
u_!=\Lan_u\ \dashv\ u^*\ \dashv\ u_*=\Ran_u.
\end{equation}
Both extension functors are fully faithful. On the ambient diagram category
define
\[
\sk_R=u_!u^*,\qquad \cosk_R=u_*u^*.
\]
The counit $\sk_R X\to X$ detects $R$-skeletal objects; the unit
$X\to\cosk_R X$ detects $R$-coskeletal objects. If $D\subseteq R$, the
skeletal and coskeletal presentations for $D$ factor through those for $R$,
by composition of Kan extensions. This is the ordering of the induced levels
used throughout the paper.

\begin{definition}\label{def:resolution}
Let $D\subseteq R\subseteq C$ be full subcategories. The level $R$
\emph{$k$-resolves $D$} if for every $A:D\to\Vk$ the canonical unit
\begin{equation}\label{eq:resolution}
\Lan_D A\longrightarrow \Ran_R\res_R\Lan_D A
\end{equation}
is an equivalence. A \emph{coefficient-relative Aufhebung}, written
$\Aup_k(D)$, is a least such full subcategory under inclusion, if it exists.
The definition for $k=\infty$ uses unrestricted spaces.
\end{definition}

The qualifier records a substantive restriction. For $k=0$, this is the
usual skeletal-to-coskeletal resolution condition on the specified ordinary
presheaf levels. For higher $k$, we extend the condition through
\eqref{eq:triple}. We do not silently identify the admissible family of
induced levels with every possible essential localization of a higher
category of diagrams.

\subsection{Matching maps and least reconstruction subsets}
\label{sec:matching}

The existence of a least full resolver over a finite poset can be
separated from the calculation of that resolver. Let $\mathcal V$ be an
infinity-category with finite limits. For $X:C\to\mathcal V$, write
\begin{equation}\label{eq:matchingsupport}
M_xX=\limop_{y\in C_{>x}}X(y),\qquad
S(X)=\{x\in C:X(x)\longrightarrow M_xX\text{ is not an equivalence}\}.
\end{equation}
The map is induced by the arrows out of $x$; at a maximal element its
target is terminal. These are matching objects for the inverse Reedy
structure obtained by assigning smaller degree to higher objects of
$C$. Matching objects and inductive extension of diagrams are standard
Reedy machinery \cite[section 3]{RiehlVerity}. The following elementary
consequence is useful because it applies to arbitrary full subposets,
not just degree cuts.

\begin{proposition}[Least reconstruction subset]\label{prop:support}
Let $C$ be a finite poset, $R\subseteq C$ full, and
$X:C\to\mathcal V$. The canonical unit
$X\to\Ran_R(X|_R)$ is an equivalence if and only if $S(X)\subseteq R$.
Consequently, for any specified family $\mathcal F$ of such diagrams,
the least full subposet from which every member is right Kan extended is
\begin{equation}\label{eq:familysupport}
S(\mathcal F)=\bigcup_{X\in\mathcal F}S(X).
\end{equation}
\end{proposition}

\begin{proof}
If $X$ is right Kan extended from $R$ and $x\notin R$, factor the
inclusion through $C\setminus\{x\}$. Composition of right Kan
extensions shows that $X$ is also right Kan extended from this
intermediate subposet. Its unit at the only omitted object is exactly
$X(x)\to M_xX$, because
$(x\downarrow(C\setminus\{x\}))=C_{>x}$. Thus $x\notin S(X)$.

Conversely, suppose $S(X)\subseteq R$. List the objects of
$C\setminus R$ in an increasing linear extension, and delete them in
that order. When $x$ is deleted, every $y>x$ is still present. The
comparison for right extension from the remaining subposet is an
equivalence at every retained object, and at $x$ it is the original
matching map $X(x)\to M_xX$. This is an equivalence by assumption.
Compose the finitely many right Kan extensions to obtain the assertion
for $R$. Applying the equivalence to every $X\in\mathcal F$ proves
the formula for a family.
\end{proof}

\begin{corollary}[Arbitrary starting levels]\label{cor:arbitraryD}
If $D\subseteq C$ is any full subposet and $\mathcal V$ also admits
the left Kan extensions in question, the least full resolver of $D$
exists and is
\begin{equation}\label{eq:arbitraryD}
D\,\cup\!\bigcup_{A\in\Fun(D,\mathcal V)}S(\Lan_D A).
\end{equation}
In particular this holds for every $\mathcal V_k$ used here.
\end{corollary}

This is an existence result, not a general algorithm: the quantifier
over $A$ can be substantial. It also locates the role of leastness in
our finite setting. Once a diagram family and its matching tests are
fixed, the least subset is already forced; the substantive problem is
to compute it and to justify the chosen family. No new general theory
of Reedy matching objects is claimed.

\begin{example}[Why the finiteness hypothesis matters]\label{ex:infinite}
Take $C=\mathbb N$ with its usual order and $D=\{0\}$.
Every strict upper interval has a least element, so the matching support
of every constant diagram is empty. Nevertheless a full $R$ containing
$0$ resolves constants exactly when it is unbounded: then each
$(n\downarrow R)$ has a least element, while for bounded $R$ an upper
comma is eventually empty. There is no inclusion-minimal unbounded
subset containing $0$, because any positive member may be deleted.
Thus the finite matching-support conclusion cannot be extended to
arbitrary infinite posets merely by using the same formula.
\end{example}

\subsection{Resolution of an initial one-object level}

Let $C$ have an initial object $b$, and let $D=\{b\}$. Since the comma
category $(D\downarrow x)$ has one object and its identity for every $x$,
left extension from $D$ is the constant diagram:
\begin{equation}\label{eq:constantlan}
\Lan_D A\simeq\Delta_C A.
\end{equation}
This holds for every coefficient category under consideration, including
arbitrary spaces. There is no homotopy-colimit defect in the left extension
of this particular starting level.

\begin{lemma}[Constant-diagram criterion]\label{lem:criterion}
Let $R\subseteq C$ be full and contain $b$. For finite $k\geq0$, the
following are equivalent:
\begin{enumerate}[label=(\roman*)]
\item $R$ $k$-resolves $D$;
\item $\tau_{\leq k}\B(x\downarrow R)\simeq *$ for every $x\in C$.
\end{enumerate}
For $k=\infty$, replace (ii) by contractibility of every
$\B(x\downarrow R)$.
\end{lemma}

\begin{proof}
Pointwise right Kan extension and \eqref{eq:constantlan} identify the
comparison at $x$ with
\begin{equation}\label{eq:mapcomparison}
A\longrightarrow\limop_{(x\downarrow R)}\Delta A
\simeq\Map\bigl(\B(x\downarrow R),A\bigr).
\end{equation}
The equivalence is the standard formula for a limit of a constant diagram of
spaces. The map sends a point of $A$ to the corresponding constant map.
For $A$ a $k$-type, the target depends only on
$\tau_{\leq k}\B(x\downarrow R)$. Thus (ii) implies (i).
Conversely, if \eqref{eq:mapcomparison} is an equivalence for every
$A\in\Vk$, the map
$\tau_{\leq k}\B(x\downarrow R)\to *$ induces equivalences on mapping
spaces into every object of $\Vk$. Yoneda implies that it is an
equivalence. The same argument without truncation proves the unrestricted
case. Equivalences of diagrams are detected pointwise.
\end{proof}

Our connectivity convention is that a nonempty space $T$ is
\emph{$k$-connected} when $\tau_{\leq k}T\simeq *$. Thus 0-connected
means connected, and 1-connected means simply connected. We state the
convention because cofinality references can shift the index by one. The
lemma is a constant-diagram instance of truncated cofinality \cite{Du};
the short proof fixes exactly the special case needed here.

For sets the criterion sees only connected components. For groupoids it also
sees the fundamental group. For unrestricted spaces it requires the whole
homotopy type of each comma nerve to be trivial. An ordinary resolution can
therefore leave an obstruction detectable at the next coefficient level.

\subsection{An intrinsic criterion for finite posets}
\label{sec:intrinsic}

For constant diagrams, the matching-map calculation specializes to a
criterion depending only on the indexing poset. It is also a form of
truncated cofinality \cite{Du}; the matching argument supplies a direct
proof without an additional cofinality theorem.

For a finite poset $C$, let $C_{>x}$ be the strict upper interval at $x$.
Define
\begin{equation}\label{eq:essentialset}
E_k(C)=\{x\in C:\tau_{\leq k}\B(C_{>x})\not\simeq *\}.
\end{equation}
For $k=\infty$, replace the condition by noncontractibility. Every
maximal element belongs to $E_k(C)$, because its upper interval is empty.

\begin{theorem}[Upper-interval criterion]\label{thm:intrinsic}
For finite $k\geq0$, a finite poset $C$, and a full subposet $R$,
\begin{equation}\label{eq:intrinsic}
\bigl[\tau_{\leq k}\B(x\downarrow R)\simeq *\text{ for every }x\in C\bigr]
\quad\Longleftrightarrow\quad E_k(C)\subseteq R.
\end{equation}
The analogous statement holds without truncation for $k=\infty$.
If $b$ is initial and $D=\{b\}$, it follows that
\begin{equation}\label{eq:generalAufhebung}
\Aup_k(D)=\{b\}\cup E_k(C)
\end{equation}
among full-subcategory levels containing $b$.
\end{theorem}

\begin{proof}
For $X=\Delta_C A$, the matching map at $x$ is
$A\to\Map(\B(C_{>x}),A)$. As $A$ ranges over all $k$-types,
Yoneda identifies the union of their matching supports with $E_k(C)$.
For unrestricted spaces the same argument uses no truncation.
Apply \cref{prop:support} to this family: right extension from $R$
recovers every constant diagram exactly when $E_k(C)\subseteq R$.
The pointwise mapping-space calculation in the proof of
\cref{lem:criterion}, which does not require an initial object,
identifies it with the comma-nerve condition in \eqref{eq:intrinsic}.
If $b$ is initial, left generation
from $D$ is constant, so \cref{cor:arbitraryD} gives the last formula.
\end{proof}

For $k=0$, the dual lower-interval description is already part of the
theory of least full initial subposets: see Dey--Lesnick
\cite[Remark 3.6]{DeyLesnick} and Br\"ustle--Desrochers--Leblanc
\cite[Theorem C and Remark 2]{BrustleEtAl}. The higher criterion is the matching-map consequence just proved,
consistent with standard truncated cofinality. The untruncated
reasoning is related to established point-deletion methods for finite
spaces \cite{BarmakMinian}. We use the intrinsic form to organize and
calculate a resolution profile, rather than claiming a new cofinality
theorem.

\section{A family with an exact resolution profile}
\label{sec:family}

\subsection{The indexing posets}

Fix $N\geq3$. Let $\Pcat_N$ consist of nonempty proper subsets of $[N]$,
ordered by inclusion. Define
\begin{equation}\label{eq:Cdef}
\Ccat_N=\{b<a\}\star\Pcat_N,
\end{equation}
where this means $b<a<S$ for every $S\in\Pcat_N$, with the original
inclusion order on $\Pcat_N$. Thus $|\Ccat_N|=2^N$. The object $b$ is
initial. We compare all full subposets containing $D=\{b\}$.

For $0\leq k\leq N-3$, define
\begin{equation}\label{eq:Rdef}
R_k=\{b\}\cup\{S\in\Pcat_N:|S|\geq N-k-2\}.
\end{equation}
The object $a$ is omitted. The phrase ``top $k+2$ ranks'' refers to
cardinalities $N-k-2,\ldots,N-1$, not to the existence of a largest
element. The proper Boolean poset has several maximal elements.

\subsection{The standard topological ingredient}

\begin{lemma}[Simplex skeleton]\label{lem:skeleton}
Let $\Delta^{m-1}$ be the simplex on $m$ vertices, and let
$1\leq q\leq m-2$. Its $q$-skeleton has homotopy type
\begin{equation}\label{eq:wedge}
\bigvee^{\binom{m-1}{q+1}}S^q.
\end{equation}
It is $(q-1)$-connected and has nonzero reduced integral homology in degree
$q$. The order complex of the poset of nonempty proper subsets of an
$m$-element set is $S^{m-2}$ for $m\geq2$.
\end{lemma}

\begin{proof}
Write $K=\sk_q\Delta^{m-1}$ and choose a vertex $v$. Its closed star is
$L=v*\sk_{q-1}\Delta^{m-2}$, a contractible subcomplex. Every simplex
of dimension less than $q$ belongs to $L$. The $q$-simplices outside $L$
are exactly those not containing $v$, of which there are
$\binom{m-1}{q+1}$. Their boundaries belong to $L$. Collapsing $L$
gives the stated wedge. Since $L\hookrightarrow K$ is a cofibration and
$L$ is nonempty and contractible, the quotient map is a homotopy
equivalence. For the final assertion, the order complex is the barycentric
subdivision of the boundary of an $(m-1)$-simplex.
\end{proof}

This is elementary standard topology. The proof makes the quantitative
calculation independent of a black-box connectivity assertion.

\begin{lemma}[The omitted-face nerves]\label{lem:nerve}
Let $0\leq k\leq N-3$. If $x\in R_k$, then
$\B(x\downarrow R_k)$ is contractible. At $x=a$ it is homotopy
equivalent to
\begin{equation}\label{eq:anervecalc}
\bigvee^{\binom{N-1}{k+2}}S^{k+1}.
\end{equation}
If $x=S\in\Pcat_N\setminus R_k$ and $m=N-|S|$, then
\begin{equation}\label{eq:facenervecalc}
\B(S\downarrow R_k)\simeq
\bigvee^{\binom{m-1}{k+2}}S^{k+1}.
\end{equation}
\end{lemma}

\begin{proof}
For $x\in R_k$, the identity arrow is initial in the comma category.
For $a$, the comma category is the poset of proper subsets of $[N]$ of
cardinality at least $N-k-2$. Complementation reverses its order and gives
the poset of nonempty subsets of $[N]$ of cardinality at most $k+2$.
Its order complex is the barycentric subdivision of
$\sk_{k+1}\Delta^{N-1}$. Reversing a poset does not change the
homeomorphism type of its order complex. Apply \cref{lem:skeleton}.

For an omitted $S$, the comma category consists of selected supersets
$T\supset S$. Taking complements in $[N]$ identifies its opposite with
the nonempty subsets of $[N]\setminus S$ of cardinality at most $k+2$.
Here $m\geq k+3$, since $S$ is omitted. Apply the same argument to
$\sk_{k+1}\Delta^{m-1}$.
\end{proof}

\subsection{The main theorem}

\begin{theorem}[Exact coefficient-relative Aufhebung]\label{thm:main}
For $D=\{b\}$ in $\Ccat_N$, the least resolving full subposet exists at
every coefficient level. It is
\begin{equation}\label{eq:main}
\Aup_k(D)=
\begin{cases}
R_k,&0\leq k\leq N-3,\\
\Ccat_N,&N-2\leq k\leq\infty.
\end{cases}
\end{equation}
In particular the levels at $k=0,1,\ldots,N-2$ form a strictly increasing
chain of $N-1$ different least resolutions.
\end{theorem}

\begin{proof}
Compute the set $E_k(\Ccat_N)$ in \cref{thm:intrinsic}.
The strict upper interval of $b$ has initial object $a$, so $b$ never
belongs to this essential set. The strict upper interval of $a$ is
$\Pcat_N$, with classifying space $S^{N-2}$. Hence $a$ belongs
exactly when $k\geq N-2$, or when $k=\infty$.

For $S\in\Pcat_N$, let $m=N-|S|$. If $m=1$, its strict upper
interval is empty, so $S$ is always essential. If $m\geq2$, the
strict upper interval is the proper nonempty Boolean poset on $m$
elements, with classifying space $S^{m-2}$. It is $k$-connected
exactly when $m-2>k$. Thus for finite $k$, the essential subsets are
precisely those with $|S|\geq N-k-2$; every subset is essential for
unrestricted coefficients. Formula \eqref{eq:generalAufhebung} gives
\eqref{eq:main}, including leastness against all full subposets.
The wedge calculation in \cref{lem:nerve} independently displays the
comma nerves of each proposed resolver.

Successive $R_k$ add the nonempty rank $|S|=N-k-3$, until
$R_{N-3}$ contains every object except $a$. The final step adds $a$.
All inclusions in the stated chain are strict.
\end{proof}

The theorem does not rely on an unproved monotonicity principle for arbitrary
full subcategories. Leastness is checked against every admissible competitor.
It also avoids confusing a discrete-input calculation with a statement about
all higher types: the proof uses every $A\in\Vk$ through
\cref{lem:criterion}.

\begin{corollary}[Unbounded length]\label{cor:unbounded}
For every $r\geq1$, some finite indexing poset and fixed initial level have
at least $r$ distinct least coefficient-relative resolutions. Moreover,
\begin{equation}\label{eq:size}
|\Aup_k(D)|=1+\sum_{j=1}^{k+2}\binom Nj\quad(0\leq k\leq N-3),
\qquad |\Aup_{N-2}(D)|=2^N.
\end{equation}
\end{corollary}

\begin{proof}
Take $N\geq\max(3,r+1)$ and apply \cref{thm:main}. The count follows
by complementation of the selected proper subsets.
\end{proof}

\subsection{A compact family and the smallest first jump}
\label{sec:compact}

The Boolean family records the ranks of explicit wedges, but exponential
size is unnecessary for a long profile. Let $Q_m$ be the ordinal sum
of $m$ two-element antichains, with layers $1,\ldots,m$ from bottom
to top. A chain selects at most one vertex in each layer, so
$\B Q_m$ is the join of $m$ copies of $S^0$, hence $S^{m-1}$.
These are the classical finite sphere models of McCord, studied in
minimal-model theory by Barmak--Minian \cite{BarmakMinimal}.

\begin{proposition}[Linear-size profiles]\label{prop:compact}
For $m\geq2$ put $C'_m=\{b<a\}\star Q_m$, with $D=\{b\}$.
For $0\leq k\leq m-2$, its least resolver consists of $b$ and the
last $k+2$ layers of $Q_m$, and has $2k+5$ objects. For
$k\geq m-1$, including $k=\infty$, the least resolver is $C'_m$.
Thus $2m+2$ objects realize $m$ distinct least levels. For $m=2$
this gives a six-object Set-to-groupoid jump, and no smaller finite
poset with an initial one-object starting level can exhibit that jump.
\end{proposition}

\begin{proof}
The strict upper interval of an element in layer $j<m$ is $Q_{m-j}$,
with nerve $S^{m-j-1}$; the top layer has empty upper intervals.
The interval above $a$ has nerve $S^{m-1}$, while the interval above
$b$ has initial object $a$. The upper-interval criterion gives exactly
the displayed levels. At every omitted object of a proper least
resolver, the selected upper comma is $Q_{k+2}$, with nerve $S^{k+1}$.

For minimality of six, a strict Set-to-groupoid change requires some
$x\ne b$ with connected but not simply connected strict upper nerve.
A connected poset on at most three objects has contractible order
complex: it is a point, a tree, or the filled triangle of a total
order. The upper interval therefore has at least four objects.
Adding the distinct objects $x$ and $b$ gives at least six.
\end{proof}

The compact calculation is an application of standard sphere models,
not a claim to have discovered those models. It shows that the
coefficient sensitivity is not an artifact of the size of a Boolean
lattice. The larger family retains its separate quantitative use:
the ranks of its wedges measure multiple simultaneous obstructions.

\subsection{An explicit witness at the next coefficient level}

\begin{proposition}[Detection of missing coherence]\label{prop:witness}
For $0\leq k\leq N-3$, the level $R_k$ fails to $(k+1)$-resolve $D$.
For $A=\K(\Z,k+1)$, the failure is visible at $a$: the target of
\eqref{eq:mapcomparison} has connected components indexed by
\begin{equation}\label{eq:componentwitness}
\Z^{\binom{N-1}{k+2}},
\end{equation}
whereas its source is connected. The constant-map comparison reaches only
the zero component.
\end{proposition}

\begin{proof}
Put $q=k+1\geq1$. By \cref{lem:nerve}, the indexing space at $a$ is
$T\simeq\bigvee^h S^q$, where $h=\binom{N-1}{k+2}>0$.
Eilenberg--Mac Lane representability gives
\[
\pi_0\Map(T,\K(\Z,q))\cong H^q(T;\Z)\cong\Z^h.
\]
The space $\K(\Z,q)$ is connected. Its constant maps have zero degree
on every sphere and belong to the zero cohomology class. Hence the
comparison is not an equivalence.
\end{proof}

This is a cohomological witness to a homotopical failure, not an assertion
that cohomology generally classifies homotopy types. The preceding wedge
identification justifies the calculation. No inference from acyclicity to
contractibility is used.

\subsection{Relation to ordinary topos levels}

At $k=0$ we are in the ordinary presheaf topos
$\Fun(\Ccat_N,\Set)=\Set^{\Ccat_N}$. Restriction to a full subposet
has the triple \eqref{eq:triple}, with left-exact middle functor, and induces
an essential subtopos. Finite posets have no nontrivial retract identifications. Apply the
finite-site rigidity result to $\Ccat_N^{\mathrm{op}}$
\cite{Marques}. Its subtoposes are induced by full subposets. More
explicitly, the topology $J_R$ declares a sieve at $x$ covering exactly
when it contains every arrow $r\to x$ with $r\in R$. The objects of
$R$ are exactly the $J_R$-irreducible objects. Thus inclusion of the
selected subsets agrees with the subtopos ordering, and a level above
$D$ must include $b$. The set-valued member is therefore an ordinary
Aufhebung calculation among all such levels. For higher coefficients
the theorem retains its explicit full-subcategory qualification.

These posets are not a counterexample to a special cohesive-site theorem.
In $\Ccat_N^{\mathrm{op}}$, the terminal object is $b$, but there is no
map $b\to x$ for $x\ne b$. The usual hypothesis that each site object
has points therefore fails. The canonical components functor also fails
to preserve binary products: representables at two distinct maximal
faces each have singleton colimit, whereas their pointwise product is
empty. Our comparison concerns the abstract resolution condition, not
a refutation of a result with those additional cohesive hypotheses.

\section{Examples: the loop and nonconstant generation}
\label{sec:example}

\subsection{The eight-object Boolean member}

Take $N=3$. The six elements of $\Pcat_3$ are the vertices and edges of
a triangle, ordered by incidence. Their nerve is the subdivided boundary
circle. The full category adds $b<a$ beneath them, as in
\cref{fig:eight}. The theorem gives
\[
\Aup_0(D)=\{b\}\cup\Pcat_3,\qquad
\Aup_1(D)=\Aup_\infty(D)=\Ccat_3.
\]

\begin{figure}[htbp]
\centering
\begin{tikzpicture}[x=1.4cm,y=1.15cm,
 dot/.style={circle,draw=ink,fill=white,inner sep=2.7pt},
 edge/.style={draw=ink!65,line width=.65pt},
 every node/.style={font=\small}]
\node[dot,label=above:{$12$}] (e12) at (-1.45,3) {};
\node[dot,label=above:{$13$}] (e13) at (0,3) {};
\node[dot,label=above:{$23$}] (e23) at (1.45,3) {};
\node[dot,label=left:{$1$}] (v1) at (-1.45,2) {};
\node[dot,label=right:{$2$}] (v2) at (0,2) {};
\node[dot,label=right:{$3$}] (v3) at (1.45,2) {};
\foreach \v/\e in {v1/e12,v1/e13,v2/e12,v2/e23,v3/e13,v3/e23}
\draw[edge] (\v)--(\e);
\node[dot,draw=accent,fill=accent!12,label=right:{$a$}] (a) at (0,.8) {};
\node[dot,fill=ink,label=below:{$b$}] (b) at (0,0) {};
\foreach \v in {v1,v2,v3}\draw[edge] (a)--(\v);
\draw[edge] (b)--(a);
\node[align=left,text=ink] at (3.1,2.5)
 {$\B(a\downarrow R_0)\simeq S^1$\\[4pt]
  $\tau_{\leq0}S^1\simeq *$\\[4pt]
  $\tau_{\leq1}S^1\not\simeq *$};
\end{tikzpicture}
\caption{The eight-object poset $\Ccat_3$. The set-valued least resolver
contains every object except $a$; the space-valued least resolver also
requires $a$. This is a Hasse diagram: arrows point upward and transitive
arrows are omitted. The six incidence edges form a circle in the order
complex, despite crossings in the drawing.}
\label{fig:eight}
\end{figure}

For a set $A$, compatible assignments of its elements over a connected
indexing category are constant. Thus
\[
\limop_{\Pcat_3}\Delta A\cong A.
\]
There is no other omitted-object test, since only $a$ is omitted. Leastness
also has a direct description. Every maximal edge object must be present:
otherwise its comma category is empty. Each vertex must then be present:
if it is absent, the two incident maximal edges give a disconnected comma
category. All six objects are forced.

For a space $A$, however,
\[
\limop_{\Pcat_3}\Delta A\simeq\Map(S^1,A).
\]
The comparison selects constant loops. Taking $A=\K(\Z,1)$ produces
components indexed by $\Z$. The earlier resolution did not require this
loop data to vanish; it did not test it. Once groupoids become admissible,
the loop becomes visible. Adding $a$ gives the relevant comma category an
initial object and repairs the failure.

\begin{table}[htbp]
\centering
\begin{tabular}{@{}crrrr@{}}
\toprule
Indexing size & Sets & Groupoids & 2-types & 3-types\\
\midrule
$|\Ccat_3|=8$  &7&8&8&8\\
$|\Ccat_4|=16$ &11&15&16&16\\
$|\Ccat_5|=32$ &16&26&31&32\\
\bottomrule
\end{tabular}
\caption{Cardinalities of the least resolving full subposets. The starting
level and ambient indexing data are fixed within each row. Every row
eventually stabilizes at the whole indexing category.}
\label{tab:profiles}
\end{table}

These numbers are not a measurement of philosophical progress. They count
objects in a constructed poset. Their role is to distinguish senses of
resolution that ordinary notation could otherwise conceal.

\subsection{A nonconstant nine-object example}\label{sec:nonconstant}

The constant generation of the main family makes its comparison unusually
transparent. The same first jump survives genuinely nonconstant generation.
Start with $p<a<\Pcat_3$, and add an object $q$ with $q<1$, hence
$q<12,13$, and no other new comparisons. Here $1,2,3$ denote singleton
faces, while $12,13,23$ denote two-element faces. Let $D=\{p,q\}$;
its two objects are incomparable. For a pair $(A,B)$ of coefficient
objects, left extension produces the diagram $X=\Lan_D(A,B)$ with
\begin{equation}\label{eq:nonconstant}
\begin{array}{c|c}
\text{positions}&\text{value of }X\\\hline
p,a,2,3,23&A\\
q&B\\
1,12,13&A\sqcup B.
\end{array}
\end{equation}
The nonidentity transition maps are identities or the appropriate
coproduct inclusions. The lower comma categories for generation are
discrete, so these coproducts agree on old inputs in every coefficient
category. The diagrams are not constant, even when $A$ and $B$ are
connected nonempty types.

\begin{proposition}\label{prop:nonconstant}
In this nine-object poset the least full resolver of $D$ is $C\setminus
\{a\}$ for sets, and $C$ for all $k\geq1$, including unrestricted
spaces.
\end{proposition}

\begin{proof}
Taking $B=\varnothing$ and $A$ a two-element set forces all six faces:
the matching map at an edge is $A\to *$, and at a vertex it is the
diagonal $A\to A\times A$. Both fail to be bijective. Thus any
resolver already contains $D\cup\Pcat_3$.

To test the only remaining position $a$, restrict $X$ to
$J=\Pcat_3$. Color its $A$ summands by $0$ and its $B$ summands by
$1$, giving a natural map $X|_J\to\Delta_J\{0,1\}$. Since
$\B J\simeq S^1$ is connected, the limit of the latter is
$\{0,1\}$. The fiber of $\lim_JX$ over $0$ is the limit of the
constant $A$ diagram, hence $\Map(S^1,A)$. Its fiber over $1$ is
empty: that fiber diagram has empty values at $2,3,23$. Limits
preserve these pullbacks. Consequently
\[
M_aX=\limop_JX\simeq\Map(S^1,A),
\]
and its matching map is the constant-map comparison from $A$.
For sets it is always an equivalence. For $k\geq1$, take
$A=\K(\Z,1)$ and $B=*$: the comparison misses all nonzero winding
components. Apply \cref{prop:support,cor:arbitraryD}.
\end{proof}

This example isolates the two operations even without constant output
from left generation: every earlier generated diagram is preserved,
while the right reconstruction detects additional coherence.
